%%% MA_00_Organisation_and_Module_Walkthrough.tex -- reorganised; frames moved from the old MA_xx files (see comments)

%%% [moved from MA_0_Organisation_and_Module_Walkthrough.tex, frame 0]
% \section{Introduction and Module Walkthrough}

\begin{frame}{Welcome!}
    \mytext{
    Welcome to the module.

    This is a module for BIM\\
    \mbox{Betriebswirtschaft / Internationales Management}. 

    Who else is here?
    }
    \mycomment{
    Make sure you fulfil the technical requirements, such as registering in CampinO and Aulis, and the necessary permissions.
    }
\end{frame}

\begin{frame}{Communication in this Module}
    \mytext{
    \begin{itemize}
        \item Course language is English only
        \item we address each other by first name \emph{in English}
        \item In class, say your names every time you speak
        \item ...and background / course
        \item state how you want to be addressed (nickname?)
        \item communicate special requirements early
        \item Main point of communication is Aulis/Ilias
        \item when using email:
        \begin{itemize}
            \item use \texttt{\makeatletter@\makeatother stud.hs-bremen.de} only
            \item do not send attachments 
            \item state your name and class            
        \end{itemize}
    \end{itemize}
    }
    \mycomment{
    English is communication language also for administrative topics.

    We lecturers are not permitted to communicate personal information to email addresses outside of hsb.
    }
    
\end{frame}

\begin{frame}{Communication in this Module}
    \showmybitmap{screenshot-memebers.png}
    \mycomment{
    \begin{itemize}
        \item complete your AULIS profile:
        \begin{itemize}
            \item upload a picture
            \item click on \enquote{profil freigeben} (different tab!)
        \end{itemize}
    \end{itemize}
    }
\end{frame}

\begin{frame}{Communication in this Module}
    \mytext{
    \begin{itemize}
        \item subsequent lectures will be Zoom
        \item live, using green screen
        \item cameras switched on 
        \item perfectly interactive!
        \item avoid joining together on one machine
        \item arrive on time and have 5 minutes for soundcheck
        \item small talk at the beginning encouraged
        \item ask questions any time
    \end{itemize}
    }
    \mycomment{
        Make sure you test your equipment. A good microphone or headsets help. Joining together impairs communication, as not all of you will be visible and audible. It also is prone to create feedback loops and disturb other people in the same room.
    }
\end{frame}



%%% [moved from MA_0_Organisation_and_Module_Walkthrough.tex, frame 5]
\begin{frame}{Portfolio Exam}
    \mytext{
    The exam is portfolio. That means it is a unitary exam, and the final grades are composed in view of all results of all parts and of  students at the end. 

    Both parts are also pass-fail. You must pass both.
    
    \begin{itemize}
        \item 30\% presentation of an economy, 20 minutes
        \item 70\% written exam at the end of term
    \end{itemize}
    }
    \mycomment{

        That means also that there will not be a publication of exact grades for any part of the portfolio.

        This also means that it is not possible to repeat a single part of the portfolio -- you can only repeat the whole exam with all parts (and the entire module) if you fail, or upon proof of illness or other force majeure.

        For all other issues regarding the exam: note that you have to deal with your Prüfungsamt and Prüfungsausschuss. The lecturer is responsible only for matters of content. The Prüfungsamt also sets the date of the written exam.
    }
\end{frame}

%%% [moved from MA_0_Organisation_and_Module_Walkthrough.tex, frame 6]
\begin{frame}{Presentation: Video}
    \mytext{
    Choose an Economy: a \textbf{country, state or region} (\enquote{Gebietskörperschaft} / jurisdiction / monetary area)!

   Zoom into an area, a province, or a city.

   Zoom into an industry, a policy, or a challenge.

    The result shall be a video of 20 minutes. 
    }
    \mycomment{
        Rather than holding your presentation live in person, I want you to present on video. We will discuss technical options. 

        A video of a presentation is a video of slides in HSB style with the presenter audible, or even visible. 
        
    }
\end{frame}


\begin{frame}{Announcing Your Topics Online}
    \mytext{
    Use the AULIS forum to announce your chosen topic, so we don't have doubles.
    }
\end{frame}


%%% [moved from MA_0_Organisation_and_Module_Walkthrough.tex, frame 7]
\begin{frame}{Presentation: Video}
    \mytext{    
    Your presentation has two parts; part one -- you present  economic data; part two -- you present the government's economic policy.
    
    \begin{itemize}
        \item 20 minutes, no fancy animations
        \item using HSB style guide / slide layout
        \item (live) recorded video of presentation
        \item students audible, visible if possible
        \item students will give feedback and grade each other
        \item[] that feeds into the grading.
    \end{itemize}
    }
    \mycomment{
        To receive a grade, you yourself must also comment and grade a minimum number of other videos.
        
        This will allow all students to scrutinise your performance and hence make grading more transparent.
    }
\end{frame}

%%% [moved from MA_0_Organisation_and_Module_Walkthrough.tex, frame 8]
\begin{frame}{Presentation: Video}
    \mytext{
    Part I: Present core economic data: 

    \begin{itemize}
        \item figures: present, history, trends
        \item compare to Germany / Bremen
    \end{itemize}

    }
    \mycomment{
    Use a \enquote{zoom-in} or \enquote{focus} technology. Present the country data as framework, and present the largest components.
    
    When looking at a large state that is organised as a federal state (USA, China, Germany, India, Brazil etc.), zoom into a sub-unit (federal unit).

    }

\end{frame}

%%% [moved from MA_0_Organisation_and_Module_Walkthrough.tex, frame 9]
\begin{frame}{Presentation: Video}
    \mytext{
    Part I: Core economic data

        \begin{itemize}
            \item core geo- and demographic data
            \item constitution, institutions
            \item GDP, per head, average wage, inequality
            \item degree of industrialisation
            \item core industries / prices
            \item inflation, currency
            \item labour markets: employment, education
            \item import / export dependency
            \item trade openness
        \end{itemize}
    }
    \mycomment{

    When discussing the industry or core industries, focus on the most important one. For instance, Germany is a highly industrialised country; the most important ones are car manufacturing, machines, and chemistry; for chemistry, these are the latest numbers... etc.

    You can already come to some conclusion at the end of this part, as to where the main challenges lie.
    }

\end{frame}

%%% [moved from MA_0_Organisation_and_Module_Walkthrough.tex, frame 10]
\begin{frame}{Presentation: Video}
    \mytext{
    Part II: The State and Economic Policy
        \begin{itemize}
            \item state levels, budget, social security budget / welfare state size
            \item stated objectives, real objectives
            \item investment subsidies
            \item production subsidies
            \item regulation and standards
            \item wages, social policy, education
            \item export subsidies
            \item tax rebates, special economic areas
            \item government procurement
            \item innovation policy
            \item monetary policy 
        \end{itemize}
    }
    \mycomment{
        In the second part, you explain the current government's economic policy. You must demonstrate critical understanding. How do you view these policies?
        
        You should also mention: What are the main opposition party proposals? What do they criticise and what do they propose instead? What does the international economic community, the IMF and the World Bank, propose?    
    }
\end{frame}

\begin{frame}{Criteria}
    \mytext{
    Presentations will be judged / graded by -
    \begin{enumerate}
        \item clarity and specificity
        \item correctness, scientific method
        \item completeness, depth
    \end{enumerate}

    \textbf{Use the style guide of Hochschule Bremen.}
    }
\end{frame}
\begin{frame}{Do's and Don'ts}
    \mycomment{
    This is a presentation about an economic situation, not a sales pitch. Provide us with everything we need to know -- and leave everything aside that we do not need to know.
    }
    \mytext{
    \begin{itemize}
        \item NO AI illustrations
        \item NO embellishments
        \item NO fancy animations or embedded videos
    \end{itemize}
    But - 
    \begin{enumerate}
        \item every student must speak 
        \item speaker be stated on lower third (slide or video)
        \item use technology of your choice (ppt, obs)
    \end{enumerate}
    }
\end{frame}

\begin{frame}{Tutorials (\enquote{Modulübung})}
    \mytext{
    Not mandatory. I show you the technology I use.
    \begin{itemize}
        \item \LaTeX\ using Overleaf
        \item OBS studio for video
        \item Tools, AI
    \end{itemize}
    I am looking forward to your hints as well!
    }
    
\end{frame}
